\documentclass[../../main.tex]{subfiles}% 注意这里的文件路径不能用 ./main.tex ,否则用latexmk编译子文件会报错
\graphicspath{{\subfix{./image/}}{\subfix{../../image/}}} % 指定图片引用路径,后续可以直接使用图片文件名
% 如果图片存放在上级目录,则使用 ../../image/ 作为备用路径

% 例如:
% \begin{figure}[H]
% \centering
% \includegraphics[width=0.8\textwidth]{图.png}
% \caption{}
% \label{figure:图}
% \end{figure}
% 注意:上述\label{}一定要放在\caption{}之后,否则引用图片序号会只会显示??.

\begin{document}

\section{外微分式和外微分}

\begin{definition}
设$E^3$为3维欧氏空间，正则曲面$S$的参数方程为$\boldsymbol{r}(u^1,u^2)$，$(u^1,u^2)\in D\subset E^2$，则
\begin{align}
\mathrm{d}\boldsymbol{r}=\boldsymbol{r}_\alpha\mathrm{d}u^\alpha.
\label{eq:dr_surface}
\end{align}
曲面$S$在点$p$的切空间$T_pS$以$\{\boldsymbol{r}_1,\boldsymbol{r}_2\}$为自然基，$\{\mathrm{d}u^1,\mathrm{d}u^2\}$为其对偶基；称$T_pS$的对偶空间为\textbf{余切空间}，记为$T_p^*S$，其元素为$S$在$p$点的\textbf{余切向量}，即切空间上的线性函数。
\end{definition}

\begin{definition}
设$D\subset E^n$为$n$维欧氏空间区域，$(x^1,\cdots,x^n)$为$D$上的直角坐标系；$U\subset E^n$为另一区域，$(u^1,\cdots,u^n)$为$U$上直角坐标系。若存在$C^\infty$微分同胚$f:U\to D,\ g:D\to U$，满足
\begin{align*}
x^\alpha=f^\alpha(u^1,\cdots,u^n),\quad u^\alpha=g^\alpha(x^1,\cdots,x^n),\quad 1\leqslant\alpha\leqslant n,
\end{align*}
则称$(u^1,\cdots,u^n)$为$D$上的\textbf{曲纹坐标系}。
区域$D$在点$p$的切空间自然基为$\left\{\frac{\partial}{\partial u^1},\cdots,\frac{\partial}{\partial u^n}\right\}$，余切空间基为$\{\mathrm{d}u^1,\cdots,\mathrm{d}u^n\}$。
\end{definition}

\begin{proposition}
$n$个$C^\infty$函数$x^\alpha=f^\alpha(u^1,\cdots,u^n)\ (1\leqslant\alpha\leqslant n)$能在点$(x_0^1,\cdots,x_0^n)$邻域内引入曲纹坐标系的充要条件为其Jacobi行列式满足
\begin{align}
\frac{\partial(f^1,\cdots,f^n)}{\partial(u^1,\cdots,u^n)}\neq0.
\label{eq:jacobi_nonzero}
\end{align}
\end{proposition}
\begin{proof}
\textbf{必要性}: 由微分同胚复合求导，
\begin{align*}
\frac{\partial g^\alpha}{\partial x^\gamma}\frac{\partial f^\gamma}{\partial u^\beta}=\delta_\beta^\alpha,
\end{align*}
故Jacobi行列式非零。

\textbf{充分性}: 由反函数定理，Jacobi行列式非零则局部存在$C^\infty$逆映射$u^\alpha=g^\alpha(x^1,\cdots,x^n)$，满足复合恒等式，故可引入曲纹坐标系.

\end{proof}

\begin{definition}
$n$维\textbf{微分流形}是由$n$维欧氏空间的局部区域通过曲纹坐标的$C^\infty$相容变换拼接而成的拓扑空间；其每一点可定义切空间、余切空间，$\{\mathrm{d}u^1,\cdots,\mathrm{d}u^n\}$为余切空间的局部基。
\end{definition}

\begin{definition}
设$D\subset E^n$为区域，$(u^1,\cdots,u^n)$为$D$上曲纹坐标系。若在每点$p\in D$指定一个$r$次外形式
\begin{align}
\varphi(p)&=\frac{1}{r!}\varphi_{i_1\cdots i_r}(u^1,\cdots,u^n)\mathrm{d}u^{i_1}\wedge\cdots\wedge\mathrm{d}u^{i_r}\nonumber\\
&=\sum\limits_{1\leqslant i_1<\cdots<i_r\leqslant n}\varphi_{i_1\cdots i_r}(u^1,\cdots,u^n)\mathrm{d}u^{i_1}\wedge\cdots\wedge\mathrm{d}u^{i_r},
\label{eq:r_differential_form}
\end{align}
其中系数$\varphi_{i_1\cdots i_r}$关于指标$i_1,\cdots,i_r$反对称且为$C^\infty$函数，则称$\varphi$为$D$上的\textbf{$r$次外微分式}。
\end{definition}

\begin{example}
$D$上的1次外微分式即1次微分式。对$C^\infty$函数$f:D\to\mathbb{R}$，其微分
\begin{align*}
\mathrm{d}f=\frac{\partial f}{\partial u^i}\mathrm{d}u^i
\end{align*}
为$D$上的1次外微分式。
\end{example}

\begin{example}
设$S$为$E^3$中正则参数曲面，参数方程$\boldsymbol{r}=\boldsymbol{r}(u^1,u^2)$，第一基本形式$\mathrm{I}=g_{\alpha\beta}\mathrm{d}u^\alpha\mathrm{d}u^\beta$。定义
\begin{align}
\mathrm{d}\sigma=\sqrt{g_{11}g_{22}-(g_{12})^2}\mathrm{d}u^1\wedge\mathrm{d}u^2,
\label{eq:area_element}
\end{align}
则$\mathrm{d}\sigma$为$S$上的2次外微分式，称为$S$的\textbf{面积元素}。
\end{example}

\begin{proposition}
曲面$S$的面积元素$\mathrm{d}\sigma$在\textbf{保定向容许参数变换}下形式不变。
\end{proposition}
\begin{proof}
设保定向参数变换$\tilde{u}^\alpha=\tilde{u}^\alpha(u^1,u^2)$满足
\begin{align*}
\frac{\partial(\tilde{u}^1,\tilde{u}^2)}{\partial(u^1,u^2)}>0,
\end{align*}
新旧参数下第一基本形式的度量矩阵满足
\begin{align*}
\begin{pmatrix}g_{11}&g_{12}\\g_{12}&g_{22}\end{pmatrix}=J\begin{pmatrix}\tilde{g}_{11}&\tilde{g}_{12}\\\tilde{g}_{12}&\tilde{g}_{22}\end{pmatrix}J^\mathrm{T},
\end{align*}
其中$J=\begin{pmatrix}\frac{\partial\tilde{u}^1}{\partial u^1}&\frac{\partial\tilde{u}^2}{\partial u^1}\\\frac{\partial\tilde{u}^1}{\partial u^2}&\frac{\partial\tilde{u}^2}{\partial u^2}\end{pmatrix}$，取行列式得
\begin{align*}
g_{11}g_{22}-(g_{12})^2=(\det J)^2(\tilde{g}_{11}\tilde{g}_{22}-(\tilde{g}_{12})^2).
\end{align*}
又外积满足
\begin{align*}
\mathrm{d}\tilde{u}^1\wedge\mathrm{d}\tilde{u}^2=\frac{\partial(\tilde{u}^1,\tilde{u}^2)}{\partial(u^1,u^2)}\mathrm{d}u^1\wedge\mathrm{d}u^2=\det J\,\mathrm{d}u^1\wedge\mathrm{d}u^2,
\end{align*}
结合保定向条件$\det J>0$，代入得
\begin{align*}
\sqrt{\tilde{g}_{11}\tilde{g}_{22}-(\tilde{g}_{12})^2}\mathrm{d}\tilde{u}^1\wedge\mathrm{d}\tilde{u}^2=\sqrt{g_{11}g_{22}-(g_{12})^2}\mathrm{d}u^1\wedge\mathrm{d}u^2,
\end{align*}
故$\mathrm{d}\sigma$形式不变.

\end{proof}

\begin{remark}
面积元素$\mathrm{d}\sigma(\boldsymbol{a},\boldsymbol{b})=\pm|\boldsymbol{a}\times\boldsymbol{b}|$，即切向量$\boldsymbol{a},\boldsymbol{b}$张成的平行四边形的\textbf{有向面积}，故$\mathrm{d}\sigma$刻画曲面的局部面积微元。
\end{remark}

\begin{definition}
设$\varphi=\frac{1}{r!}\varphi_{i_1\cdots i_r}\mathrm{d}u^{i_1}\wedge\cdots\wedge\mathrm{d}u^{i_r}$为$D$上$r$次外微分式，定义其\textbf{外微分}$\mathrm{d}\varphi$为$r+1$次外微分式：
\begin{align}
\mathrm{d}\varphi&=\frac{1}{r!}\mathrm{d}\varphi_{i_1\cdots i_r}\wedge\mathrm{d}u^{i_1}\wedge\cdots\wedge\mathrm{d}u^{i_r}\nonumber\\
&=\frac{1}{r!}\frac{\partial\varphi_{i_1\cdots i_r}}{\partial u^j}\mathrm{d}u^j\wedge\mathrm{d}u^{i_1}\wedge\cdots\wedge\mathrm{d}u^{i_r}.
\label{eq:exterior_derivative}
\end{align}
特别地，零次外微分式（光滑函数）的外微分即为普通微分。
\end{definition}

\begin{example}
设$(u,v,w)$为$E^3$中曲纹坐标系，取
\begin{align*}
\omega&=\alpha(u,v,w)\mathrm{d}u+\beta(u,v,w)\mathrm{d}v+\gamma(u,v,w)\mathrm{d}w,\\
\eta&=P(u,v,w)\mathrm{d}v\wedge\mathrm{d}w+Q(u,v,w)\mathrm{d}w\wedge\mathrm{d}u+R(u,v,w)\mathrm{d}u\wedge\mathrm{d}v,
\end{align*}
其中$\alpha,\beta,\gamma,P,Q,R\in C^\infty$，则
\begin{align*}
\mathrm{d}\omega&=\left(\frac{\partial\gamma}{\partial v}-\frac{\partial\beta}{\partial w}\right)\mathrm{d}v\wedge\mathrm{d}w+\left(\frac{\partial\alpha}{\partial w}-\frac{\partial\gamma}{\partial u}\right)\mathrm{d}w\wedge\mathrm{d}u+\left(\frac{\partial\beta}{\partial u}-\frac{\partial\alpha}{\partial v}\right)\mathrm{d}u\wedge\mathrm{d}v,\\
\mathrm{d}\eta&=\left(\frac{\partial P}{\partial u}+\frac{\partial Q}{\partial v}+\frac{\partial R}{\partial w}\right)\mathrm{d}u\wedge\mathrm{d}v\wedge\mathrm{d}w.
\end{align*}
\end{example}

\begin{theorem}
外微分算子$\mathrm{d}$满足：
\begin{enumerate}[(1)]
\item 线性性：$\mathrm{d}(\varphi^1+\varphi^2)=\mathrm{d}\varphi^1+\mathrm{d}\varphi^2,\quad \mathrm{d}(c\varphi^1)=c\mathrm{d}\varphi^1,\ \forall c\in\mathbb{R}$；
\item 幂零性：$\mathrm{d}\circ\mathrm{d}=0$，即$\mathrm{d}(\mathrm{d}\varphi)=0$；
\item Leibniz法则：若$\varphi$为$r$次外微分式，则$\mathrm{d}(\varphi\wedge\psi)=\mathrm{d}\varphi\wedge\psi+(-1)^r\varphi\wedge\mathrm{d}\psi$。
\end{enumerate}
\end{theorem}
\begin{proof}
(1) 由外微分的定义直接可得。

(2) 设$\varphi=\frac{1}{r!}\varphi_{i_1\cdots i_r}\mathrm{d}u^{i_1}\wedge\cdots\wedge\mathrm{d}u^{i_r}$，则
\begin{align*}
\mathrm{d}\varphi&=\frac{1}{r!}\frac{\partial\varphi_{i_1\cdots i_r}}{\partial u^j}\mathrm{d}u^j\wedge\mathrm{d}u^{i_1}\wedge\cdots\wedge\mathrm{d}u^{i_r},\\
\mathrm{d}(\mathrm{d}\varphi)&=\frac{1}{r!}\frac{\partial^2\varphi_{i_1\cdots i_r}}{\partial u^k\partial u^j}\mathrm{d}u^k\wedge\mathrm{d}u^j\wedge\mathrm{d}u^{i_1}\wedge\cdots\wedge\mathrm{d}u^{i_r}.
\end{align*}
因$C^\infty$函数二阶偏导可交换，即$\frac{\partial^2\varphi_{i_1\cdots i_r}}{\partial u^k\partial u^j}=\frac{\partial^2\varphi_{i_1\cdots i_r}}{\partial u^j\partial u^k}$，而$\mathrm{d}u^k\wedge\mathrm{d}u^j=-\mathrm{d}u^j\wedge\mathrm{d}u^k$，故求和后项两两抵消，$\mathrm{d}(\mathrm{d}\varphi)=0$。

(3) 由线性性，仅需验证单项式$\varphi=a\mathrm{d}u^{i_1}\wedge\cdots\wedge\mathrm{d}u^{i_r},\ \psi=b\mathrm{d}u^{j_1}\wedge\cdots\wedge\mathrm{d}u^{j_s}$的情形：
\begin{align*}
\mathrm{d}(\varphi\wedge\psi)&=\mathrm{d}(ab)\wedge\mathrm{d}u^{i_1}\wedge\cdots\wedge\mathrm{d}u^{i_r}\wedge\mathrm{d}u^{j_1}\wedge\cdots\wedge\mathrm{d}u^{j_s}\\
&=(a\mathrm{d}b+b\mathrm{d}a)\wedge\mathrm{d}u^{i_1}\wedge\cdots\wedge\mathrm{d}u^{i_r}\wedge\mathrm{d}u^{j_1}\wedge\cdots\wedge\mathrm{d}u^{j_s}\\
&=\mathrm{d}\varphi\wedge\psi+(-1)^r\varphi\wedge\mathrm{d}\psi.
\end{align*}

\end{proof}

\begin{corollary}
\begin{enumerate}[(1)]
\item 若$f$为零次外微分式（光滑函数），则$\mathrm{d}(f\psi)=\mathrm{d}f\wedge\psi+f\mathrm{d}\psi$；
\item 若$\varphi$为1次外微分式，则$\mathrm{d}(\varphi\wedge\psi)=\mathrm{d}\varphi\wedge\psi-\varphi\wedge\mathrm{d}\psi$。
\end{enumerate}
\end{corollary}
\begin{proof}
(1) 零次形式$r=0$，代入Leibniz法则得$\mathrm{d}(f\psi)=\mathrm{d}f\wedge\psi+(-1)^0f\mathrm{d}\psi=\mathrm{d}f\wedge\psi+f\mathrm{d}\psi$。

(2) 一次形式$r=1$，代入得$\mathrm{d}(\varphi\wedge\psi)=\mathrm{d}\varphi\wedge\psi+(-1)^1\varphi\wedge\mathrm{d}\psi=\mathrm{d}\varphi\wedge\psi-\varphi\wedge\mathrm{d}\psi$.

\end{proof}

\begin{example}
在欧氏空间$E^3$的笛卡儿直角坐标系下，将向量场的场论公式与外微分式的二次外微分为零的性质进行对照。
\end{example}
\begin{solution}

设$(u,v,w)$为$E^3$的笛卡儿直角坐标系，$\omega=(\alpha,\beta,\gamma)$，$\eta=(P,Q,R)$为区域$D$上的连续可微向量场，则旋度与散度为
\begin{align*}
\mathrm{rot}(\alpha,\beta,\gamma) &= \left(\frac{\partial\gamma}{\partial v}-\frac{\partial\beta}{\partial w},\frac{\partial\alpha}{\partial w}-\frac{\partial\gamma}{\partial u},\frac{\partial\beta}{\partial u}-\frac{\partial\alpha}{\partial v}\right), \\
\mathrm{div}(P,Q,R) &= \frac{\partial P}{\partial u}+\frac{\partial Q}{\partial v}+\frac{\partial R}{\partial w}.
\end{align*}
场论中满足恒等式$\mathrm{div}\circ\mathrm{rot}=0$，代入向量场$\omega$得
\begin{align*}
\mathrm{div}\big(\mathrm{rot}(\alpha,\beta,\gamma)\big) &= \mathrm{div}\left(\frac{\partial\gamma}{\partial v}-\frac{\partial\beta}{\partial w},\frac{\partial\alpha}{\partial w}-\frac{\partial\gamma}{\partial u},\frac{\partial\beta}{\partial u}-\frac{\partial\alpha}{\partial v}\right) \\
&= \frac{\partial}{\partial u}\left(\frac{\partial\gamma}{\partial v}-\frac{\partial\beta}{\partial w}\right)+\frac{\partial}{\partial v}\left(\frac{\partial\alpha}{\partial w}-\frac{\partial\gamma}{\partial u}\right)+\frac{\partial}{\partial w}\left(\frac{\partial\beta}{\partial u}-\frac{\partial\alpha}{\partial v}\right)=0.
\end{align*}

将向量场对应为外微分式：
\begin{enumerate}[(1)]
\item 向量场$\omega=(\alpha,\beta,\gamma)$对应1次外微分式$\tilde{\omega}=\alpha\mathrm{d}u+\beta\mathrm{d}v+\gamma\mathrm{d}w$，其外微分为
\begin{align*}
\mathrm{d}\tilde{\omega} &= \left(\frac{\partial\gamma}{\partial v}-\frac{\partial\beta}{\partial w}\right)\mathrm{d}v\land\mathrm{d}w+\left(\frac{\partial\alpha}{\partial w}-\frac{\partial\gamma}{\partial u}\right)\mathrm{d}w\land\mathrm{d}u+\left(\frac{\partial\beta}{\partial u}-\frac{\partial\alpha}{\partial v}\right)\mathrm{d}u\land\mathrm{d}v,
\end{align*}
即$\omega$的旋度对应$\tilde{\omega}$的外微分。
\item 向量场$\eta=(P,Q,R)$对应2次外微分式$\tilde{\eta}=P\mathrm{d}v\land\mathrm{d}w+Q\mathrm{d}w\land\mathrm{d}u+R\mathrm{d}u\land\mathrm{d}v$，其外微分为
\begin{align*}
\mathrm{d}\tilde{\eta} &= \left(\frac{\partial P}{\partial u}+\frac{\partial Q}{\partial v}+\frac{\partial R}{\partial w}\right)\mathrm{d}u\land\mathrm{d}v\land\mathrm{d}w,
\end{align*}
即$\eta$的散度对应$\tilde{\eta}$的外微分。
\end{enumerate}

因此场论公式$\mathrm{div}\circ\mathrm{rot}\,\omega=0$等价于$\mathrm{d}(\mathrm{d}\tilde{\omega})=0$，即
\begin{align*}
\mathrm{d}(\mathrm{d}\tilde{\omega}) = \left[\frac{\partial}{\partial u}\left(\frac{\partial\gamma}{\partial v}-\frac{\partial\beta}{\partial w}\right)+\frac{\partial}{\partial v}\left(\frac{\partial\alpha}{\partial w}-\frac{\partial\gamma}{\partial u}\right)+\frac{\partial}{\partial w}\left(\frac{\partial\beta}{\partial u}-\frac{\partial\alpha}{\partial v}\right)\right]\mathrm{d}u\land\mathrm{d}v\land\mathrm{d}w=0.
\end{align*}
该结论不依赖于直角坐标系，对任意曲纹坐标系、任意维数空间的任意次外微分式均成立。

\end{solution}

\begin{theorem}
设$\varphi$是定义在$n$维区域$D$上的$r$次外微分式，在曲纹坐标系$(u^1,\cdots,u^n)$下表示为
\begin{align*}
\varphi = \frac{1}{r!}\varphi_{i_1\cdots i_r}(u^1,\cdots,u^n)\mathrm{d}u^{i_1}\land\cdots\land\mathrm{d}u^{i_r},
\end{align*}
在另一曲纹坐标系$(\tilde{u}^1,\cdots,\tilde{u}^n)$下表示为
\begin{align*}
\varphi = \frac{1}{r!}\tilde{\varphi}_{i_1\cdots i_r}(\tilde{u}^1,\cdots,\tilde{u}^n)\mathrm{d}\tilde{u}^{i_1}\land\cdots\land\mathrm{d}\tilde{u}^{i_r},
\end{align*}
其中$\varphi_{i_1\cdots i_r}$与$\tilde{\varphi}_{i_1\cdots i_r}$对下指标均反对称，则满足
\begin{align}
\mathrm{d}\varphi_{i_1\cdots i_r}\land\mathrm{d}u^{i_1}\land\cdots\land\mathrm{d}u^{i_r} = \mathrm{d}\tilde{\varphi}_{i_1\cdots i_r}\land\mathrm{d}\tilde{u}^{i_1}\land\cdots\land\mathrm{d}\tilde{u}^{i_r}. \label{eq:exterior-form-invariance}
\end{align}
\end{theorem}
\begin{proof}
设两坐标系间的容许坐标变换为$\tilde{u}^i=\tilde{u}^i(u^1,\cdots,u^n)\ (1\leqslant i\leqslant n)$，则
\begin{align}
\mathrm{d}\tilde{u}^i = \frac{\partial\tilde{u}^i}{\partial u^j}\mathrm{d}u^j. \label{eq:coord-diff}
\end{align}
代入外微分式得
\begin{align*}
\varphi_{i_1\cdots i_r}\mathrm{d}u^{i_1}\land\cdots\land\mathrm{d}u^{i_r} &= \tilde{\varphi}_{j_1\cdots j_r}\mathrm{d}\tilde{u}^{j_1}\land\cdots\land\mathrm{d}\tilde{u}^{j_r} \\
&= \tilde{\varphi}_{j_1\cdots j_r}\frac{\partial\tilde{u}^{j_1}}{\partial u^{i_1}}\cdots\frac{\partial\tilde{u}^{j_r}}{\partial u^{i_r}}\mathrm{d}u^{i_1}\land\cdots\land\mathrm{d}u^{i_r}.
\end{align*}
由系数的反对称性，比较两端得
\begin{align}
\varphi_{i_1\cdots i_r} = \tilde{\varphi}_{j_1\cdots j_r}\frac{\partial\tilde{u}^{j_1}}{\partial u^{i_1}}\cdots\frac{\partial\tilde{u}^{j_r}}{\partial u^{i_r}}. \label{eq:coeff-transform}
\end{align}
对\eqref{eq:coeff-transform}求微分
\begin{align*}
\mathrm{d}\varphi_{i_1\cdots i_r} &= \frac{\partial}{\partial u^k}\left(\tilde{\varphi}_{j_1\cdots j_r}\frac{\partial\tilde{u}^{j_1}}{\partial u^{i_1}}\cdots\frac{\partial\tilde{u}^{j_r}}{\partial u^{i_r}}\right)\mathrm{d}u^k \\
&= \left(\frac{\partial\tilde{\varphi}_{j_1\cdots j_r}}{\partial\tilde{u}^l}\frac{\partial\tilde{u}^l}{\partial u^k}\frac{\partial\tilde{u}^{j_1}}{\partial u^{i_1}}\cdots\frac{\partial\tilde{u}^{j_r}}{\partial u^{i_r}}+\tilde{\varphi}_{j_1\cdots j_r}\frac{\partial^2\tilde{u}^{j_1}}{\partial u^k\partial u^{i_1}}\cdots\frac{\partial\tilde{u}^{j_r}}{\partial u^{i_r}}+\cdots+\tilde{\varphi}_{j_1\cdots j_r}\frac{\partial\tilde{u}^{j_1}}{\partial u^{i_1}}\cdots\frac{\partial^2\tilde{u}^{j_r}}{\partial u^k\partial u^{i_r}}\right)\mathrm{d}u^k.
\end{align*}
因此
\begin{align*}
\mathrm{d}\varphi_{i_1\cdots i_r}\land\mathrm{d}u^{i_1}\land\cdots\land\mathrm{d}u^{i_r} &= \left(\frac{\partial\tilde{\varphi}_{j_1\cdots j_r}}{\partial\tilde{u}^l}\frac{\partial\tilde{u}^l}{\partial u^k}\frac{\partial\tilde{u}^{j_1}}{\partial u^{i_1}}\cdots\frac{\partial\tilde{u}^{j_r}}{\partial u^{i_r}}+\sum\text{高阶偏导项}\right)\mathrm{d}u^k\land\mathrm{d}u^{i_1}\land\cdots\land\mathrm{d}u^{i_r}.
\end{align*}
高阶混合偏导项满足
\begin{align*}
\frac{\partial^2\tilde{u}^{j_1}}{\partial u^k\partial u^{i_1}}\cdots\frac{\partial\tilde{u}^{j_r}}{\partial u^{i_r}}\mathrm{d}u^k\land\mathrm{d}u^{i_1}\land\cdots\land\mathrm{d}u^{i_r} &= \frac{1}{2}\left(\frac{\partial^2\tilde{u}^{j_1}}{\partial u^k\partial u^{i_1}}-\frac{\partial^2\tilde{u}^{j_1}}{\partial u^{i_1}\partial u^k}\right)\cdots\frac{\partial\tilde{u}^{j_r}}{\partial u^{i_r}}\mathrm{d}u^k\land\mathrm{d}u^{i_1}\land\cdots\land\mathrm{d}u^{i_r}=0,
\end{align*}
故仅保留第一项：
\begin{align*}
\mathrm{d}\varphi_{i_1\cdots i_r}\land\mathrm{d}u^{i_1}\land\cdots\land\mathrm{d}u^{i_r} &= \frac{\partial\tilde{\varphi}_{j_1\cdots j_r}}{\partial\tilde{u}^l}\frac{\partial\tilde{u}^l}{\partial u^k}\frac{\partial\tilde{u}^{j_1}}{\partial u^{i_1}}\cdots\frac{\partial\tilde{u}^{j_r}}{\partial u^{i_r}}\mathrm{d}u^k\land\mathrm{d}u^{i_1}\land\cdots\land\mathrm{d}u^{i_r} \\
&= \frac{\partial\tilde{\varphi}_{j_1\cdots j_r}}{\partial\tilde{u}^l}\mathrm{d}\tilde{u}^l\land\mathrm{d}\tilde{u}^{j_1}\land\cdots\land\mathrm{d}\tilde{u}^{j_r} \\
&= \mathrm{d}\tilde{\varphi}_{j_1\cdots j_r}\land\mathrm{d}\tilde{u}^{j_1}\land\cdots\land\mathrm{d}\tilde{u}^{j_r}.
\end{align*}

\end{proof}
\begin{remark}
外微分是定义在光滑流形上的整体算子：流形上每点邻域存在局部坐标系，不同局部坐标系间为容许坐标变换，而外微分算子与局部坐标系选取无关，因此可将流形上的$r$次外微分式映射为$r+1$次外微分式。
\end{remark}

\begin{definition}
设$\sigma:D\to\tilde{D}$为从$n$维区域$D$到$m$维区域$\tilde{D}$的连续可微映射，其坐标表示为
\begin{align}
\tilde{u}^\alpha = \tilde{u}^\alpha(u^1,\cdots,u^n),\quad \alpha=1,\cdots,m. \label{eq:diff-map}
\end{align}
对$\tilde{D}$上的$r$次外微分式
\begin{align*}
\tilde{\varphi} = \frac{1}{r!}\sum\limits_{1\leqslant\alpha_1,\cdots,\alpha_r\leqslant m}\tilde{\varphi}_{\alpha_1\cdots\alpha_r}(\tilde{u}^1,\cdots,\tilde{u}^m)\mathrm{d}\tilde{u}^{\alpha_1}\land\cdots\land\mathrm{d}\tilde{u}^{\alpha_r},
\end{align*}
定义$\tilde{\varphi}$通过$\sigma$在$D$上的\textbf{拉回}$\sigma^*\tilde{\varphi}$为将\eqref{eq:diff-map}代入$\tilde{\varphi}$得到的$D$上的$r$次外微分式：
\begin{align}
\sigma^*\tilde{\varphi} &= \frac{1}{r!}\tilde{\varphi}_{\alpha_1\cdots\alpha_r}\big(\tilde{u}^1(u^1,\cdots,u^n),\cdots,\tilde{u}^m(u^1,\cdots,u^n)\big)\frac{\partial\tilde{u}^{\alpha_1}}{\partial u^{i_1}}\cdots\frac{\partial\tilde{u}^{\alpha_r}}{\partial u^{i_r}}\mathrm{d}u^{i_1}\land\cdots\land\mathrm{d}u^{i_r}, \label{eq:pullback-def}
\end{align}
式中采用Einstein和式约定。
\end{definition}

\begin{theorem}
设$\sigma:D\to\tilde{D}$为连续可微映射，对$\tilde{D}$上任意外微分式$\varphi,\psi$，拉回映射$\sigma^*$满足：
\begin{enumerate}[(1)]
\item $\sigma^*(\varphi+\psi)=\sigma^*\varphi+\sigma^*\psi$；
\item $\sigma^*(\varphi\land\psi)=\sigma^*\varphi\land\sigma^*\psi$；
\item $\sigma^*(\mathrm{d}\varphi)=\mathrm{d}(\sigma^*\varphi)$.
\end{enumerate}
\end{theorem}
\begin{proof}
性质(1)(2)可由拉回映射的定义直接验证；性质(3)的证明与定理2.2一致，定理2.2是性质(3)在恒等坐标变换下的特殊情形。

\end{proof}

\begin{theorem}
设$G$为$n$维欧氏空间$E^n$中$r$维有向曲面上的区域，$\partial G$为$G$的诱导定向边界，$\omega$为$G$上的$r-1$次外微分式，则
\begin{align}
\int_{\partial G}\omega = \int_G\mathrm{d}\omega. \label{eq:stokes-general}
\end{align}
\end{theorem}
\begin{remark}
式\eqref{eq:stokes-general}称为\textbf{Stokes公式}，是微积分中积分公式的统一形式，其证明可参考微分几何相关文献。
\end{remark}

\begin{example}
在二维、三维欧氏空间中，将Stokes公式\eqref{eq:stokes-general}具体表述为经典积分公式。
\end{example}
\begin{solution}
\begin{enumerate}[(1)]
\item 当$n=r=2$，令$\omega=P(x,y)\mathrm{d}x+Q(x,y)\mathrm{d}y$，则
\begin{align*}
\mathrm{d}\omega = \left(\frac{\partial Q}{\partial x}-\frac{\partial P}{\partial y}\right)\mathrm{d}x\land\mathrm{d}y,
\end{align*}
代入\eqref{eq:stokes-general}得
\begin{align*}
\int_{\partial G}P\mathrm{d}x+Q\mathrm{d}y = \iint_G\left(\frac{\partial Q}{\partial x}-\frac{\partial P}{\partial y}\right)\mathrm{d}x\land\mathrm{d}y,
\end{align*}
即经典的\textbf{Green公式}。

\item 当$n=3,r=2$，令$\omega=P\mathrm{d}x+Q\mathrm{d}y+R\mathrm{d}z$，则
\begin{align*}
\mathrm{d}\omega &= \left(\frac{\partial R}{\partial y}-\frac{\partial Q}{\partial z}\right)\mathrm{d}y\land\mathrm{d}z+\left(\frac{\partial P}{\partial z}-\frac{\partial R}{\partial x}\right)\mathrm{d}z\land\mathrm{d}x+\left(\frac{\partial Q}{\partial x}-\frac{\partial P}{\partial y}\right)\mathrm{d}x\land\mathrm{d}y,
\end{align*}
代入\eqref{eq:stokes-general}得
\begin{align*}
\int_{\partial G}P\mathrm{d}x+Q\mathrm{d}y+R\mathrm{d}z &= \iint_G\left(\frac{\partial R}{\partial y}-\frac{\partial Q}{\partial z}\right)\mathrm{d}y\land\mathrm{d}z+\left(\frac{\partial P}{\partial z}-\frac{\partial R}{\partial x}\right)\mathrm{d}z\land\mathrm{d}x+\left(\frac{\partial Q}{\partial x}-\frac{\partial P}{\partial y}\right)\mathrm{d}x\land\mathrm{d}y,
\end{align*}
即经典的\textbf{Stokes公式}。

\item 当$n=r=3$，代入\eqref{eq:stokes-general}得
\begin{align*}
\iint_{\partial G}P\mathrm{d}y\land\mathrm{d}z+Q\mathrm{d}z\land\mathrm{d}x+R\mathrm{d}x\land\mathrm{d}y = \iiint_G\left(\frac{\partial P}{\partial x}+\frac{\partial Q}{\partial y}+\frac{\partial R}{\partial z}\right)\mathrm{d}x\land\mathrm{d}y\land\mathrm{d}z,
\end{align*}
即经典的\textbf{Gauss公式}。
\end{enumerate}
\end{solution}

\begin{example}
设$\varphi = yz\mathrm{d}x + \mathrm{d}z$,$\psi = \sin z\mathrm{d}x + \cos z\mathrm{d}y$,$\eta = \mathrm{d}y + z\mathrm{d}z$,计算:
\begin{enumerate}[(1)]
\item $\varphi \wedge \psi$,$\psi \wedge \eta$,$\eta \wedge \varphi$;
\item $\mathrm{d}\varphi$,$\mathrm{d}\psi$,$\mathrm{d}\eta$.
\end{enumerate}
\end{example}

\begin{example}
设$f,g$是两个光滑函数，利用外微分的运算法则将下列各式化简:
\begin{enumerate}[(1)]
\item $\mathrm{d}(f\mathrm{d}g + g\mathrm{d}f)$;
\item $\mathrm{d}((f-g)(\mathrm{d}f + \mathrm{d}g))$;
\item $\mathrm{d}((f\mathrm{d}g) \wedge (g\mathrm{d}f))$;
\item $\mathrm{d}(g\mathrm{d}f) + \mathrm{d}(f\mathrm{d}g)$.
\end{enumerate}
\end{example}

\begin{example}
假定$x,y,z$是$u,v,w$的光滑函数，证明:
\begin{align*}
\mathrm{d}x \wedge \mathrm{d}y \wedge \mathrm{d}z = \frac{\partial(x,y,z)}{\partial(u,v,w)}\mathrm{d}u \wedge \mathrm{d}v \wedge \mathrm{d}w.
\end{align*}
\end{example}

\begin{example}
设
\begin{align*}
\varphi = \frac{1}{r!}\varphi_{i_1\cdots i_r}\mathrm{d}u^{i_1} \wedge \cdots \wedge \mathrm{d}u^{i_r}, \\
\mathrm{d}\varphi = \frac{1}{(r+1)!}\alpha_{i_1\cdots i_{r+1}}\mathrm{d}u^{i_1} \wedge \cdots \wedge \mathrm{d}u^{i_{r+1}},
\end{align*}
其中系数函数$\varphi_{i_1\cdots i_r}$,$\alpha_{i_1\cdots i_{r+1}}$关于下指标都是反对称的.证明:
\begin{align*}
\alpha_{i_1\cdots i_{r+1}} = \sum\limits_{\beta=1}^{r+1}(-1)^{\beta+1}\frac{\partial \varphi_{i_1\cdots \hat{i}_\beta\cdots i_{r+1}}}{\partial u^{i_\beta}}.
\end{align*}
在这里，记号“$\hat{\ }$”表示去掉在它下方的指标字母.
\end{example}

\begin{example}
设$(r,\varphi,\theta)$是$E^3$中的球坐标系，即
\begin{align*}
x = r\cos\theta\cos\varphi,\quad y = r\cos\theta\sin\varphi,\quad z = r\sin\theta.
\end{align*}
\begin{enumerate}[(1)]
\item 求相应的自然标架场;
\item 将$\mathrm{d}x \wedge \mathrm{d}y + \mathrm{d}y \wedge \mathrm{d}z + \mathrm{d}z \wedge \mathrm{d}x$,$\mathrm{d}x \wedge \mathrm{d}y \wedge \mathrm{d}z$用球坐标系表示出来.
\end{enumerate}
\end{example}

\begin{example}
设$(r,\theta,t)$是$E^3$中的柱坐标系，即
\begin{align*}
x = r\cos\theta,\quad y = r\sin\theta,\quad z = t.
\end{align*}
\begin{enumerate}[(1)]
\item 求相应的自然标架场;
\item 将$\mathrm{d}x \wedge \mathrm{d}y + \mathrm{d}y \wedge \mathrm{d}z + \mathrm{d}z \wedge \mathrm{d}x$,$\mathrm{d}x \wedge \mathrm{d}y \wedge \mathrm{d}z$用柱坐标系表示出来.
\end{enumerate}
\end{example}

\begin{example}
设$D$是$uv$平面，命映射$\sigma:D \to E^3$定义如下:
\begin{align*}
x = \frac{2u}{1+u^2+v^2},\quad y = \frac{2v}{1+u^2+v^2},\quad z = \frac{1-u^2-v^2}{1+u^2+v^2}.
\end{align*}
假定
\begin{align*}
\omega = x\mathrm{d}x + y\mathrm{d}y + z\mathrm{d}z, \\
\xi = x\mathrm{d}y \wedge \mathrm{d}z + y\mathrm{d}z \wedge \mathrm{d}x + z\mathrm{d}x \wedge \mathrm{d}y,
\end{align*}
求:
\begin{enumerate}[(1)]
\item $\sigma^*\omega$;
\item $\sigma^*\xi$;
\item $\sigma^*(\omega \wedge \xi)$;
\item $\sigma^*(\mathrm{d}\omega)$;
\item $\sigma^*(\mathrm{d}\xi)$.
\end{enumerate}
\end{example}




\end{document}